\ifdefined\gkver\else\def\gkver{student}\fi
\documentclass[\gkver]{gaokaozhenti}
\begin{document}

\chapter{2003年全国}

\begin{enumerate}
\item
%% number: 1
%% paperName: 2003年全国高考第 1 题 4 分
%% typeId: 1
%% type: 单选题
%% chapter: 原子和原子核
%% point: 原子核的组成
%% method: 无
%% score: 4
%% degree: 650
%% duplicateId: 0
%% body:
下列说法中正确的是\xzanswer{A}
\fourchoices[answer=A]
{质子与中子的质量不等，但质量数相等}
{两个质子之间，不管距离如何，核力总是大于库仑力}
{同一种元素的原子核有相同的质量数，但中子数可以不同}
{除万有引力外，两个中子之间不存在其它相互作用力}
%% answer: A
%% memo:
\memoanswer{
质子与中子的质量数均为 1，但质量不同，故 A 错误；核力是一种短程力，随距离而变，并不总大于库仑力，故 B 错误；同位素之间中子数不同，故 C 错误；核子之间有万有引力，还有核力、库仑力，故 D 错误。
}

\item
%% number: 2
%% paperName: 2003年全国高考第 2 题 4 分
%% typeId: 2
%% type: 多选题
%% chapter: 光学
%% point: 光电效应
%% method: 无
%% score: 4
%% degree: 650
%% duplicateId: 0
%% body:
用某种单色光照射某种金属表面，发生光电效应。现将该单色光的光强减弱，则\xzanswer{AC}
\fourchoices[answer=AC]
{光电子的最大初动能不变}
{光电子的最大初动能减少}
{单位时间内产生的光电子数减少}
{可能不发生光电效应}
%% answer: AC
%% memo:
\memoanswer{
由光电效应规律知，光电子的最大初动能决定于照射光的频率，与光强无关，故 A 正确、B 错误；减弱光的强度时，单位时间内产生的光电子数减少，故 C 正确；光电效应与光强无关，仅与照射光的频率有关，故 D 错误。
}

\item
%% number: 3
%% paperName: 2003年全国高考第 3 题 4 分
%% typeId: 2
%% type: 多选题
%% chapter: 气体性质
%% point: 分子力
%% method: 无
%% score: 4
%% degree: 450
%% duplicateId: 0
%% body:
如图，甲分子固定在坐标原点 O，乙分子位于 $x$ 轴上，甲分子对乙分子的作用力与两分子间距离的关系如图中曲线所示。$F>0$ 为斥力，$F<0$ 为引力。a、b、c、d 为 $x$ 轴上四个特定的位置。现把乙分子从 a 处由静止释放，则\xzanswer{BC}
\onepicture[w=5.0cm]{\includesvg{figs/03}}
\fourchoices[answer=BC]
{乙分子从 a 到 b 做加速运动，由 b 到 c 做减速运动}
{乙分子从 a 到 c 做加速运动，到达 c 时速度最大}
{乙分子由 a 到 b 的过程中，两分子间的分子势能一直减少}
{乙分子由 b 到 d 的过程中，两分子间的分子势能一直增加}
%% answer: BC
%% memo:
\memoanswer{
乙分子从 a 到 b，再到 c 的过程，分子之间均表现为引力，显然乙分子始终做加速运动，且到达 c 点时速度最大，故 A 错误、B 正确；乙分子从 a 到 b 的过程，分子的引力一直做正功，因此两分子间的分子势能一直减少，故 C 正确；乙分子由 b 到 c 过程，分子引力仍然做正功，故两分子间的分子势能仍在减少，从 c 到 d 的过程，分子间的斥力做负功，则分子间的势能增加，故 D 错误。
}

\item
%% number: 4
%% paperName: 2003年全国高考第 4 题 4 分
%% typeId: 1
%% type: 单选题
%% chapter: 原子和原子核
%% point: 天然放射性
%% method: 无
%% score: 4
%% degree: 650
%% duplicateId: 0
%% body:
铀裂变的产物之一氪 90（$\ce{^{90}_{36}Kr}$）是不稳定的，它经过一系列衰变最终成为稳定的锆 90（$\ce{^{90}_{40}Zr}$），这些衰变是\xzanswer{B}
\fourchoices[answer=B]
{1 次 $\alpha$ 衰变，6 次 $\beta$ 衰变}
{4 次 $\beta$ 衰变}
{2 次 $\alpha$ 衰变}
{2 次 $\alpha$ 衰变，2 次 $\beta$ 衰变}
%% answer: B
%% memo:
\memoanswer{
由 $\ce{^{90}_{36}Kr}$ 衰变到 $\ce{^{90}_{40}Zr}$ 可知，质量数未变，核电荷数增加 4，故发生 4 次 $\beta$ 衰变，故 B 正确。
}

\item
%% number: 5
%% paperName: 2003年全国高考第 5 题 4 分
%% typeId: 2
%% type: 多选题
%% chapter: 电路
%% point: 电容
%% method: 无
%% score: 4
%% degree: 450
%% duplicateId: 0
%% body:
两块大小、形状完全相同的金属平板平行放置，构成以平行板电容器，与它相连接的电路如图所示，接通开关 K，电源即给电容器充电\xzanswer{BC}
\onepicture[w=6.0cm]{\includesvg{figs/05}}
\fourchoices[answer=BC]
{保持 K 接通，减小两极板间的距离，则两极板间电场的电场强度减小}
{保持 K 接通，在两极板间插入一块介质，则极板上的电量增大}
{断开 K，减小两极板间的距离，则两极板间的电势差减小}
{断开 K，在两极板间插入一块介质，则极板上的电势差增大}
%% answer: BC
%% memo:
\memoanswer{
平行板电容器的电容为 $C$，由定义式知 $C=\frac{Q}{U}$，由决定式知 $C=\frac{\varepsilon S}{4\pi kd}$。当 K 接通时，$U$ 不变，$E=\frac{U}{d}$，减小板间距，$E$ 增大，故 A 错误；在两极板间插入介质，电容 $C$ 增大，又 $U$ 不变，所以 $Q$ 增大，故 B 正确；当 K 断开时，$Q$ 不变，减小板间距，电容 $C$ 将增大，所以 $U$ 减小，故 C 正确；在两极板间插入介质，电容 $C$ 增大，所以 $U$ 减小，故 D 错误。
}

\item
%% number: 6
%% paperName: 2003年全国高考第 6 题 4 分
%% typeId: 2
%% type: 多选题
%% chapter: 气体性质
%% point: 热力学第一定律
%% method: 无
%% score: 4
%% degree: 450
%% duplicateId: 0
%% body:
一定质量的理想气体\xzanswer{CD}
\fourchoices[answer=CD]
{先等压膨胀，再等容降温，其温度必低于其始温度}
{先等温膨胀，再等压压缩，其体积必小于起始体积}
{先等容升温，再等压压缩，其温度有可能等于起始温度}
{先等容加热，再绝热压缩，其内能必大于起始内能}
%% answer: CD
%% memo:
\memoanswer{
由 $\frac{pV}{T}=C$（恒量）得 $T=\frac{pV}{C}$，先等压膨胀，体积 $V$ 将增大，再等容降温，则压强 $p$ 又减小，但 $pV$ 的值难以确定其是否增减，故 A 错误；同理，$V=\frac{T}{p}\cdot C$，等温膨胀时，压强 $p$ 减小，等压压缩时，温度 $T$ 又减小，则难以判定 $\frac{T}{p}$ 的值是否减小或增大，故 B 错误；同理，$T=\frac{pV}{C}$，先等容升温，压强 $p$ 增大，但后来等压压缩 $V$ 将减小，则 $pV$ 值可能不变，即 $T$ 可能等于起始温度，故 C 正确；先等容加热，再绝热压缩，气体的温度始终升高，则内能必定增大，故 D 正确。
}

\item
%% number: 7
%% paperName: 2003年全国高考第 7 题 4 分
%% typeId: 2
%% type: 多选题
%% chapter: 机械振动
%% point: 振动图象
%% method: 无
%% score: 4
%% degree: 450
%% duplicateId: 0
%% body:
一弹簧振子沿 $x$ 轴振动，振幅为 $4\Ucm$。振子的平衡位置位于 $x$ 轴上的 0 点。图 1 中的 a、b、c、d 为四个不同的振动状态：黑点表示振子的位置，黑点上的箭头表示运动的方向。图 2 给出的①②③④四条振动图线，可用于表示振子的振动图像\xzanswer{AD}
\twopicture[numstyle=arabic, numA=1, numB=2, labelA=2003全国07a, labelB=2003全国07b, wA=5.5cm, wB=7.5cm]{figs/07a.png}{figs/07b.png}
\fourchoices[answer=AD]
{若规定状态 a 时 $t=0$ 则图象为①}
{若规定状态 b 时 $t=0$ 则图象为②}
{若规定状态 c 时 $t=0$ 则图象为③}
{若规定状态 d 时 $t=0$ 则图象为④}
%% answer: AD
%% memo:
\memoanswer{
$t=0$ 时，a 点位移为 $3\Ucm$ 且向正方向运动，图象①，$t=0$ 时，位移为 $3\Ucm$ 且向正方向运动，故 A 正确；$t=0$ 时，b 点位移为 $2\Ucm$，而图象②，$t=0$ 时，位移为 $3\Ucm$，故 B 错误；$t=0$ 时，c 点位移为 $-2\Ucm$，速度方向向左，而图象③，$t=0$ 时，位移为 $-2\Ucm$，但速度方向向上，即向右，故 C 错误；$t=0$ 时，d 点位移为 $-4\Ucm$，速度为 0，图象④，$t=0$ 时，位移为 $-4\Ucm$，为最低点，速度为 0，故 D 正确。
}

\item
%% number: 8
%% paperName: 2003年全国高考第 8 题 3 分
%% typeId: 2
%% type: 多选题
%% chapter: 光学
%% point: 光的折射
%% method: 无
%% score: 3
%% degree: 450
%% duplicateId: 0
%% body:
如图，一玻璃柱体的横截面为半圆形。细的单色光束从空气向柱体的 O 点（半圆的圆心），产生反射光束 1 和透射光束 2。已知玻璃折射率为 $\sqrt{3}$，入射角为 $45^{\circ}$（相应的折射角为 $24^{\circ}$）。现保持入射光不变，将半圆柱绕通过 O 点垂直于图面的轴线顺时针转过 $15^{\circ}$，如图中虚线所示，则\xzanswer{BC}
\onepicture[w=4.5cm]{\includesvg{figs/08}}
\fourchoices[answer=BC]
{光束 1 转过 $15^{\circ}$}
{光束 1 转过 $30^{\circ}$}
{光束 2 转过的角度小于 $15^{\circ}$}
{光束 2 转过的角度大于 $15^{\circ}$}
%% answer: BC
%% memo:
\memoanswer{
转动前，光束 1（反射光）与入射光线间的夹角为 $A=45^{\circ}\times2=90^{\circ}$，光束 2（折射光）与入射光线间的夹角为 $B=45^{\circ}+(180^{\circ}-24^{\circ})=201^{\circ}$；转动后，反射光与入射光的夹角 $A'=60^{\circ}\times2=120^{\circ}$，由折射定律 $\frac{\sin60^{\circ}}{\sin r}=\sqrt{3}$ 得 $r=30^{\circ}$，则折射光与入射光的夹角为 $B'=60^{\circ}+(180^{\circ}-30^{\circ})=210^{\circ}$，因 $\Delta A=A'-A=30^{\circ}$，$\Delta B=B'-B=9^{\circ}$，故 BC 正确。
}

\item
%% number: 9
%% paperName: 2003年全国高考第 9 题 4 分
%% typeId: 1
%% type: 单选题
%% chapter: 原子和原子核
%% point: 玻尔模型
%% method: 无
%% score: 4
%% degree: 450
%% duplicateId: 0
%% body:
原子从一个能级跃迁到一个较低的能级时，有可能不发射光子。例如在某种条件下，铬原子的 $n=2$ 能级上的电子跃迁到 $n=1$ 能级上时并不发射光子，而是将相应的能量转交给 $n=4$ 能级上的电子，使之脱离原子，这一现象叫做俄歇效应。以这种方式脱离了原子的电子叫做俄歇电子。已知铬原子的能级公式可简化表示为 $E_{n}=-\frac{A}{n^{2}}$，式中 $n=1,2,3\ldots$ 表示不同的能级，$A$ 是正的已知常数。上述俄歇电子的动能是\xzanswer{C}
\fourchoices[answer=C]
{$\frac{3}{16}A$}
{$\frac{7}{16}A$}
{$\frac{11}{16}A$}
{$\frac{13}{16}A$}
%% answer: C
%% memo:
\memoanswer{
由题意知 $E_{2}=-\frac{A}{2^{2}}=-\frac{A}{4}$，$E_{1}=-A$，$E_{4}=-\frac{A}{16}$，又 $\Delta E=E_{2}-E_{1}=\frac{3A}{4}$，则 $E_{k}=\Delta E+E_{4}=\frac{11A}{16}$，故 C 正确。
}

\item
%% number: 10
%% paperName: 2003年全国高考第 10 题 4 分
%% typeId: 1
%% type: 单选题
%% chapter: 光学
%% point: 光的折射
%% method: 无
%% score: 4
%% degree: 450
%% duplicateId: 0
%% body:
如图，a 和 b 都是厚度均匀的平玻璃板，它们之间的夹角为 $\varphi$。一细光束以入射角 $\theta$ 从 P 点射入，$\theta>\varphi$。已知此光束由红光和蓝光组成。则当光束透过 b 板后\xzanswer{D}
\onepicture[w=6.0cm]{\includesvg{figs/10}}
\fourchoices[answer=D]
{传播方向相对于入射光方向向左偏转 $\varphi$ 角}
{传播方向相对于入射光方向向右偏转 $\varphi$ 角}
{红光在蓝光的左边}
{红光在蓝光的右边}
%% answer: D
%% memo:
\memoanswer{
光束过厚度均匀的平玻璃后，透射光与原入射光方向平行，故 A、B 错误；由于 $f_{\text{红}}<f_{\text{蓝}}$，所以 $n_{\text{红}}<n_{\text{蓝}}$，红光平移量小，故 C 错误、D 正确。
}

\item
%% number: 11
%% paperName: 2003年全国高考第 11 题 6 分
%% typeId: 4
%% type: 实验题
%% chapter: 电场
%% point: 示波器
%% method: 无
%% score: 6
%% degree: 650
%% duplicateId: 0
%% body:
图中为示波器面板，屏上显示的是一亮度很低、线条较粗且模糊不清的波形。
\begin{enumerate}
	\item
	若要增大显示波形的亮度，应调节\tkanswer[1.6cm]{辉度}旋钮。
	\item
	若要屏上波形线条变细且边缘清晰，应调节\tkanswer[1.4cm]{聚焦}旋钮。
	\item
	若要将波形曲线调至屏中央，应调节\tkanswer[1.8cm]{垂直位移}与\tkanswer[1.8cm]{水平位移}旋钮。
\end{enumerate}
\onepicture[h=7.0cm, align=right]{figs/11.png}
%% answer: 
\jdanswer{
\begin{enumerate}
	\item
	辉度
	\item
	聚焦
	\item
	垂直位移，水平位移
\end{enumerate}
}
%% memo:
\memoanswer{
本题考查示波器的使用及实验能力。有的考生没有用过甚至没有见过示波器，此题虽然容易，也无法选出正确的答案，因此物理教学要重视实验。
}

\item
%% number: 12
%% paperName: 2003年全国高考第 12 题 7 分
%% typeId: 4
%% type: 实验题
%% chapter: 直线运动
%% point: 打点计时器公式
%% method: 无
%% score: 7
%% degree: 450
%% duplicateId: 0
%% body:
实验装置如图 1 所示：一木块放在水平长木板上，左侧栓有一细软线，跨过固定在木板边缘的滑轮与一重物相连。木块右侧与打点计时器的纸带相连。在重物牵引下，木块在木板上向左运动，重物落地后，木块继续向左做匀减速运动，图 2 给出了重物落地后，打点计时器在纸带上打出的一些点，试根据给出的数据，求木块与木板间的摩擦因数 $\mu$。要求写出主要的运算过程。结果保留 2 位有效数字。（打点计时器所用交流电频率为 $50\UHz$，不计纸带与木块间的拉力。取重力加速度 $g=10\Umsq$）
\twopicture[numstyle=arabic, numA=1, numB=2, labelA=2003全国12a, labelB=2003全国12b, wA=5.0cm, wB=7.5cm, align=right]{figs/12a.png}{figs/12b.png}
%% answer: 
\jdanswer{
$\mu=0.30$
}
%% memo:
\memoanswer{
由给出的数据可知，重物落地后，木块在连续相等的时间 $T$ 内的位移分别是 $s_{1}=7.72\Ucm$，$s_{2}=7.21\Ucm$，$s_{3}=6.71\Ucm$，$s_{4}=6.25\Ucm$，$s_{5}=5.76\Ucm$，$s_{6}=5.29\Ucm$，$s_{7}=4.81\Ucm$，$s_{8}=4.31\Ucm$，以 $a$ 表示加速度，根据匀变速直线运动的规律得 $\Delta s=\frac{1}{4}[(s_{5}-s_{1})+(s_{6}-s_{2})+(s_{7}-s_{3})+(s_{8}-s_{4})]=4aT^{2}$，又知 $T=0.04\Us$，解得 $a=-3.0\Umsq$。
重物落地后木块只受摩擦力的作用，以 $m$ 表示木块的质量，由牛顿定律得 $-\mu mg=ma$，解得 $\mu=0.30$。
}

\item
%% number: 13
%% paperName: 2003年全国高考第 13 题 8 分
%% typeId: 4
%% type: 实验题
%% chapter: 电路
%% point: 闭合电路欧姆定律
%% method: 无
%% score: 8
%% degree: 450
%% duplicateId: 0
%% body:
要测量一块多用电表直流 $10\UmA$ 档的内阻 $R_{A}$（约 $40\UO$）。除此多用电表外，还有下列器材：直流电源一个（电动势 $E$ 约为 $1.5\UV$，内阻可忽略不计），电阻一个（阻值 $R$ 约为 $150\UO$），电键一个，导线若干。要求：
\begin{enumerate}
	\item
	写出实验步骤。
	\item
	给出 $R_{A}$ 的表达式。
\end{enumerate}
%% answer: 
\jdanswer{
\begin{enumerate}
	\item
	实验步骤：①用多用电表的直流电压档测量电源电动势 $E$。②用多用电表的欧姆档测电阻阻值 $R$。③将多用电表置于电流 $10\UmA$ 档，与电阻 $R$ 及电键串联后接在电源两端。合上电键，记下多用电表读数 $I$。
	\item
	$R_{A}=\frac{E}{I}-R$
\end{enumerate}
}
%% memo:
\memoanswer{
（1）实验步骤：①用多用电表的直流电压挡测量电源电动势 $E$；②用多用电表的欧姆挡测电阻阻值 $R$；③将多用电表置于直流 $10\UmA$ 挡，与电阻 $R$ 及开关串联后接在电源两端；合上开关，记下多用电表读数 $I$。

（2）$R_{A}=\frac{E}{I}-R$。
}

\item
%% number: 14
%% paperName: 2003年全国高考第 14 题 12 分
%% typeId: 6
%% type: 计算题
%% chapter: 曲线运动
%% point: 天体运动
%% method: 无
%% score: 12
%% degree: 650
%% duplicateId: 0
%% body:
据美联社 2002 年 10 月 7 日报道，天文学家在太阳系的 9 大行星之外，又发现了一颗比地球小得多的新行星，而且还测得它绕太阳公转的周期为 288 年。若把它和地球绕太阳公转的轨道都看作圆，问它与太阳的距离约是地球与太阳距离的多少倍？（最后结果可用根式表示）
%% answer: 
\jdanswer{
44（或 $\sqrt[3]{288^{2}}$）
}
%% memo:
\memoanswer{
设太阳的质量为 $M$，地球的质量为 $m_{0}$，其绕太阳公转的周期为 $T_{0}$，与太阳的距离为 $R_{0}$，公转角速度为 $\omega_{0}$；新行星的质量为 $m$，绕太阳公转的周期为 $T$，与太阳的距离为 $R$，公转角速度为 $\omega$。由万有引力定律和牛顿定律得
$G\frac{Mm}{R^{2}}=m\omega^{2}R$，$G\frac{Mm_{0}}{R_{0}^{2}}=m_{0}\omega_{0}^{2}R_{0}$，
$T=\frac{2\pi}{\omega}$，$T_{0}=\frac{2\pi}{\omega_{0}}$，
由以上各式得 $\frac{R}{R_{0}}=\left(\frac{T}{T_{0}}\right)^{\frac{2}{3}}$。
已知 $T=288$ 年，$T_{0}=1$ 年，代入解得 $\frac{R}{R_{0}}=44$（或 $\sqrt[3]{288^{2}}$）。
}

\item
%% number: 15
%% paperName: 2003年全国高考第 15 题 12 分
%% typeId: 6
%% type: 计算题
%% chapter: 运动和力
%% point: 变加速直线运动
%% method: 无
%% score: 12
%% degree: 450
%% duplicateId: 0
%% body:
当物体从高空下落时，空气阻力随速度的增大而增大，因此经过一段距离后将匀速下落，这个速度称为此物体下落的终极速度。已知球形物体速度不大时所受的空气阻力正比于速度 $v$ 且正比于球半径 $r$，即阻力 $f=krv$，$k$ 是比例系数。对于常温下的空气，比例系数 $k=3.4\times10^{-4}\UNs/\Umq$。已知水的密度 $\rho=1.0\times10^{3}\Ukgmc$，取重力加速度 $g=10\Umsq$。试求半径 $r=0.10\Umm$ 的球形雨滴在无风情况下的终极速度 $v_{r}$。（结果取两位数字）
%% answer: 
\jdanswer{
$v_{r}=1.2\Ums$
}
%% memo:
\memoanswer{
雨滴下落时受两个力作用：重力，方向竖直向下；空气阻力，方向竖直向上。当雨滴达到终极速度 $v_{T}$ 后，加速度为零，用 $m$ 表示雨滴质量，由平衡条件得 $mg-krv_{T}=0$①，
又 $m=\frac{4}{3}\pi r^{3}\rho$②，
由①②得终极速度为 $v_{T}=\frac{4\pi r^{2}\rho g}{3k}$，
代入数值得 $v_{T}=1.2\Ums$。
}

\item
%% number: 16
%% paperName: 2003年全国高考第 16 题 13 分
%% typeId: 6
%% type: 计算题
%% chapter: 电路
%% point: 电容
%% method: 无
%% score: 13
%% degree: 650
%% duplicateId: 0
%% body:
在如图所示的电路中，电源的电动势 $E=3.0\UV$，内阻 $r=1.0\UO$；电阻 $R_{1}=10\UO$，$R_{2}=10\UO$，$R_{3}=30\UO$，$R_{4}=35\UO$；电容器的电容 $C=100\UuF$。电容器原来不带电。求接通电键 K 后流过 $R_{4}$ 的总电量。
\onepicture[w=6.0cm, align=right]{\includesvg{figs/16}}
%% answer: 
\jdanswer{
$Q=2.0\times10^{-4}\UC$
}
%% memo:
\memoanswer{
由电阻的串并联公式得闭合电路的总电阻为 $R=\frac{R_{1}(R_{2}+R_{3})}{R_{1}+R_{2}+R_{3}}+r$，
由欧姆定律得，通过电源的电流为 $I=\frac{E}{R}$，
电源的路端电压为 $U=E-Ir$，
电阻 $R_{3}$ 两端的电压为 $U'=\frac{R_{3}}{R_{2}+R_{3}}U$，
通过 $R_{4}$ 的总电量就是电容器的电量为 $Q=CU'$，
代入数据解得 $Q=2.0\times10^{-4}\UC$。
}

\item
%% number: 17
%% paperName: 2003年全国高考第 17 题 13 分
%% typeId: 6
%% type: 计算题
%% chapter: 磁场
%% point: 带电粒子在组合场中的运动
%% method: 无
%% score: 13
%% degree: 650
%% duplicateId: 0
%% body:
串列加速器是用来产生高能离子的装置。图中虚线框内为其主体的原理示意图，其中加速管的中部 b 处有很高的正电势 $U$，a、c 两端均有电极接地（电势为零）。现将速度很低的负一价碳离子从 a 端输入，当离子到达 b 处时，可被设在 b 处的特殊装置将其电子剥离，成为 $n$ 价正离子，而不改变其速度大小。这些正 $n$ 价离子从 c 端飞出后进入一与其速度方向垂直的、磁感应强度为 $B$ 的匀强磁场中，在磁场中做半径为 $R$ 的圆周运动。已知碳离子的质量为 $m=2.0\times10^{-26}\Ukg$，$U=7.5\times10^{5}\UV$，$B=0.50\UT$，$n=2$，基元电荷 $e=1.6\times10^{-19}\UC$，求 $R$。
\onepicture[w=7.5cm, align=right]{\includesvg{figs/17}}
%% answer: 
\jdanswer{
$R=0.75\Um$
}
%% memo:
\memoanswer{
设碳离子到达 b 处时的速度为 $v_{1}$，从 c 端射出时的速度为 $v_{2}$，由功能关系得
$eU=\frac{1}{2}mv_{1}^{2}$①，
$\frac{1}{2}mv_{1}^{2}+neU=\frac{1}{2}mv_{2}^{2}$②，
碳离子进入磁场后，由洛伦兹力提供向心力得
$nev_{2}B=m\frac{v_{2}^{2}}{R}$③，
由以上三式可得 $R=\frac{1}{Bn}\sqrt{\frac{2mU(n+1)}{e}}$④，
由④式及题给数值可解得 $R=0.75\Um$。
}

\item
%% number: 18
%% paperName: 2003年全国高考第 18 题 13 分
%% typeId: 6
%% type: 计算题
%% chapter: 电磁感应
%% point: 电磁感应受力分析
%% method: 无
%% score: 13
%% degree: 450
%% duplicateId: 0
%% body:
如图所示，两根平行金属导轨固定在水平桌面上，每根导轨每米的电阻为 $r_{0}=0.10\UOpm$，导轨的端点 P、Q 用电阻可以忽略的导线相连，两导轨间的距离 $l=0.20\Um$。有随时间变化的匀强磁场垂直于桌面，已知磁感应强度 $B$ 与时间 $t$ 的关系为 $B=kt$，比例系数 $k=0.020\UTs$。一电阻不计的金属杆可在导轨上无摩擦低滑动，在滑动过程中保持与导轨垂直。在 $t=0$ 时刻，金属杆紧靠在 P、Q 端，在外力作用下，杆以恒定的加速度从静止开始向导轨的另一端滑动，求在 $t=6.0\Us$ 时金属杆所受的安培力。
\onepicture[w=7.0cm, align=right]{\includesvg{figs/18}}
%% answer: 
\jdanswer{
$F=1.44\times10^{-3}\UN$
}
%% memo:
\memoanswer{
以 $a$ 表示金属杆运动的加速度，在 $t$ 时刻，金属杆与初始位置的距离为 $L=\frac{1}{2}at^{2}$①，
此时杆的速度为 $v=at$②，
这时，杆与导轨构成的回路的面积为 $S=Ll$③，
导体切割磁感线产生的电动势为 $E_{1}=BLv$，
磁场变化引起磁通量变化产生的电动势为 $E_{2}=S\frac{\Delta B}{\Delta t}$，
回路中的总电动势是二者之和，有 $E=S\frac{\Delta B}{\Delta t}+BLv$④，
由题意知 $B=kt$⑤，则 $\frac{\Delta B}{\Delta t}=\frac{B(t+\Delta t)-Bt}{\Delta t}=k$⑥，
回路的总电阻为 $R=2Lr_{0}$⑦，
回路中的感应电流为 $i=\frac{E}{R}$⑧，
作用于杆的安培力的大小为 $F=Bli$⑨，
联立解得 $F=\frac{3}{2}\cdot\frac{k^{2}l^{2}}{r_{0}}t$⑩，
代入数据得 $F=1.44\times10^{-3}\UN$。
}

\item
%% number: 19
%% paperName: 2003年全国高考第 19 题 13 分
%% typeId: 6
%% type: 计算题
%% chapter: 运动和力
%% point: 动量和能量
%% method: 无
%% score: 13
%% degree: 450
%% duplicateId: 0
%% body:
图 1 所示为一根竖直悬挂的不可伸长的轻绳，下端栓一小物块 A，上端固定在 C 点且与一能测量绳的拉力的测力传感器相连。已知有一质量为 $m_{0}$ 的子弹 B 沿水平方向以速度 $v_{0}$ 射入 A 内（未穿透），接着两者一起绕 C 点在竖直面内做圆周运动。在各种阻力都可忽略的条件下测力传感器测得绳的拉力 $F$ 随时间 $t$ 的变化关系如图 2 所示。已知子弹射入的时间极短，且图 2 中 $t=0$ 为 A、B 开始以相同速度运动的时刻。根据力学规律和题中（包括图）提供的信息，对反映悬挂系统本身性质的物理量（例如 A 的质量）及 A、B 一起运动过程中的守恒量，你能求得哪些定量的结果？
\twopicture[numstyle=arabic, numA=1, numB=2, labelA=2003全国19a, labelB=2003全国19b, hA=4.4cm, hB=4.4cm, align=right]{figs/19a.png}{figs/19b.png}
%% answer: 
\jdanswer{
见解析。
}
%% memo:
\memoanswer{
由图 2 可直接看出，A、B 一起做周期性运动，运动的周期为 $T=2t_{0}$①，
令 $m$ 表示 A 的质量，$l$ 表示绳长；$v_{1}$ 表示 B 陷入 A 内时，即 $t=0$ 时 A、B 的速度（即圆周运动最低点的速度），$v_{2}$ 表示运动到最高点时的速度，$F_{1}$ 表示运动到最低点时绳的拉力，$F_{2}$ 表示运动到最高点时绳的拉力，由动量守恒得
$m_{0}v_{0}=(m_{0}+m)v_{1}$②，
在最低点和最高点处，由牛顿第二定律得
$F_{1}-(m+m_{0})g=(m+m_{0})\frac{v_{1}^{2}}{l}$③，
$F_{2}+(m+m_{0})g=(m+m_{0})\frac{v_{2}^{2}}{l}$④，
由机械能守恒得
$(m+m_{0})g\cdot2l=\frac{1}{2}(m+m_{0})v_{1}^{2}-\frac{1}{2}(m+m_{0})v_{2}^{2}$⑤，
由题图 2 可知 $F_{2}=0$⑥，
$F_{1}=F_{m}$⑦，
由以上各式可解得，反映系统性质的物理量是
$m=\frac{F_{m}}{6g}-m_{0}$⑧，
$l=\frac{36m_{0}^{2}v_{0}^{2}g}{5F_{m}^{2}}$⑨，
A、B 一起运动过程中的守恒量是机械能 $E$，若以最低点为势能的零点，则
$E=\frac{1}{2}(m+m_{0})v_{1}^{2}$⑩，
由②⑧⑩式解得 $E=\frac{3m_{0}^{2}v_{0}^{2}g}{F_{m}}$\textcircled{\scriptsize 11}。
}

\item
%% number: 20
%% paperName: 2003年全国高考第 20 题 13 分
%% typeId: 6
%% type: 计算题
%% chapter: 运动和力
%% point: 动量和能量
%% method: 无
%% score: 13
%% degree: 250
%% duplicateId: 0
%% body:
\begin{enumerate}
	\item
	如图 1，在光滑水平长直轨道上，放着一个静止的弹簧振子，它由一轻弹簧两端各联结一个小球构成，两小球质量相等。现突然给左端小球一个向右的速度 $u_{0}$，求弹簧第一次恢复到自然长度时，每个小球的速度。
	\item
	如图 2，将 $N$ 个这样的振子放在该轨道上。最左边的振子 1 被压缩至弹簧为某一长度后锁定，静止在适当位置上，这时它的弹性势能为 $E_{0}$。其余各振子间都有一定的距离。现解除对振子 1 的锁定，任其自由运动，当它第一次恢复到自然长度时，刚好与振子 2 碰撞，此后，继续发生一系列碰撞，每个振子被碰后刚好都是在弹簧第一次恢复到自然长度时与下一个振子相碰。求所有可能的碰撞都发生后，每个振子弹性势能的最大值。已知本题中两球发生碰撞时，速度交换，即一球碰后的速度等于另一球碰前的速度。
\end{enumerate}
\twopicture[numstyle=arabic, numA=1, numB=2, labelA=2003全国20a, labelB=2003全国20b, wA=7.0cm, wB=7.0cm, align=right]{figs/20a.png}{figs/20b.png}
%% answer: 
\jdanswer{
\begin{enumerate}
	\item
	$u_{1}=0$，$u_{2}=u_{0}$
	\item
	$\frac{1}{4}E_{0}$
\end{enumerate}
}
%% memo:
\memoanswer{
（1）设每个小球质量为 $m$，以 $u_{1}$、$u_{2}$ 分别表示弹簧恢复到自然长度时左右两端小球的速度，由动量守恒和能量守恒得，以向右为速度正方向，
$mu_{0}=mu_{1}+mu_{2}$①，
$\frac{1}{2}mu_{0}^{2}=\frac{1}{2}mu_{1}^{2}+\frac{1}{2}mu_{2}^{2}$②，
解得 $u_{1}=u_{0}$，$u_{2}=0$ 或 $u_{1}=0$，$u_{2}=u_{0}$。
由于振子从初始状态到弹簧恢复到自然长度的过程中，弹簧一直是压缩状态，弹性力使左端小球持续减速，使右端小球持续加速，因此应该取解 $u_{1}=0$，$u_{2}=u_{0}$③。

（2）以 $v_{1}$、$v_{1}'$ 分别表示振子 1 解除锁定后弹簧恢复到自然长度时左右两小球的速度，规定向右为速度的正方向，由动量守恒和能量守恒得
$mv_{1}+mv_{1}'=0$④，
$E_{0}=\frac{1}{2}mv_{1}^{2}+\frac{1}{2}mv_{1}'^{2}$⑤，
解得 $v_{1}=\sqrt{\frac{E_{0}}{m}}$，$v_{1}'=-\sqrt{\frac{E_{0}}{m}}$ 或 $v_{1}=-\sqrt{\frac{E_{0}}{m}}$，$v_{1}'=\sqrt{\frac{E_{0}}{m}}$。
在这一过程中，弹簧一直是压缩状态，弹性力使左端小球向左加速，右端小球向右加速，故应取解 $v_{1}=-\sqrt{\frac{E_{0}}{m}}$，$v_{1}'=\sqrt{\frac{E_{0}}{m}}$⑥。
振子 1 与振子 2 碰撞后，由于交换速度，振子 1 右端小球速度变为 0，左端小球速度仍为 $v_{1}$，此后两小球都向左运动，当它们向左的速度相同时，弹簧被拉伸至最长，弹性势能最大，设此速度为 $v_{10}$，由动量守恒得 $mv_{1}=2mv_{10}$⑦，
用 $E_{1}$ 表示最大弹性势能，由能量守恒得 $\frac{1}{2}mv_{10}^{2}+\frac{1}{2}mv_{10}^{2}+E_{1}=\frac{1}{2}mv_{1}^{2}$⑧，
解得 $E_{1}=\frac{1}{4}E_{0}$⑨。
振子 2 被碰撞后瞬间，左端小球速度为 $\sqrt{\frac{E_{0}}{m}}$，右端小球速度为 0，以后弹簧被压缩，当弹簧再恢复到自然长度时，根据（1）题结果，左端小球速度 $v_{2}=0$，右端小球速度 $v_{2}'=\sqrt{\frac{E_{0}}{m}}$，与振子 3 碰撞，由于交换速度，振子 2 右端小球速度变为 0，振子 2 静止，弹簧为自然长度，弹性势能为 $E_{2}=0$，同样分析可得 $E_{2}=E_{3}=\cdots=E_{N-1}=0$⑩。
振子 $N$ 被碰撞后瞬间，左端小球速度 $v_{N-1}'=\sqrt{\frac{E_{0}}{m}}$，右端小球速度为 0，弹簧处于自然长度，此后两小球都向右运动，弹簧被压缩，当它们向右的速度相同时，弹簧被压缩至最短，弹性势能最大，设此速度为 $v_{N0}$，由动量守恒与能量守恒得 $mv_{N-1}'=2mv_{N0}$\textcircled{\scriptsize 11}，
$\frac{1}{2}mv_{N0}^{2}+\frac{1}{2}mv_{N0}^{2}+E_{N}=\frac{1}{2}mv_{N-1}'^{2}$\textcircled{\scriptsize 12}，
解得 $E_{N}=\frac{1}{4}E_{0}$\textcircled{\scriptsize 13}。
综上可知，每个弹簧振子最大的弹性势能均为 $\frac{1}{4}E_{0}$。
}

\end{enumerate}

\end{document}
